\chapter{Construções geométricas}\label{ch:g5:geometry}

Régua, esquadro, compasso: com essas três ferramentas e um pouco de
método, qualquer figura descrita em palavras pode ser construída com
exatidão. Este capítulo constrói \omterm{def:g4:geometry:def}{perpendiculares}, \omterm{def:g4:geometry:def}{paralelas},
triângulos e os \omterm{def:g4:geometry:polygon}{quadriláteros} especiais --- primeiro o ofício, depois
a teoria (\cref{ch:g6:lines,ch:g6:shapes}).

\section{Perpendiculares e paralelas por um ponto}

\begin{method}[As duas construções básicas]\label{met:g5:geometry:basic}
\emph{Perpendicular a $d$ por $A$}: deslize o esquadro ao longo de
$d$ até a outra borda do \omterm{def:g3:shapes:rightangle}{ângulo reto} alcançar $A$; trace; marque o
\omterm{def:g3:shapes:rightangle}{ângulo reto}.

\emph{Paralela a $d$ por $A$}: encoste uma régua no esquadro, servindo
de trilho apoiado em $d$; deslize o esquadro até $A$; trace.
(Conferência: a distância entre as duas retas é a mesma em toda
parte.)
\end{method}

\begin{example}[Encadear construções]\label{ex:g5:geometry:chain}
Muitas figuras são cadeias desses dois movimentos. Um retângulo
$ABCD$ com $AB = 5$ cm e $BC = 3$ cm: trace $[AB]$; perpendicular em
$A$; perpendicular em $B$; marque $D$ e $C$ a $3$ cm do mesmo lado;
ligue. Cada passo é uma construção básica.
\end{example}

\section{Triângulos com o compasso}

\begin{method}[Triângulo a partir de três lados]\label{met:g5:geometry:triangle}
Para construir um triângulo de lados $6$ cm, $5$ cm e $4$ cm:
\begin{enumerate}
  \item trace o lado mais longo, $[AB]$ com $AB = 6$ cm;
  \item arco de centro $A$ e \omterm{def:g3:shapes:shapes}{raio} $5$ cm;
  \item arco de centro $B$ e \omterm{def:g3:shapes:shapes}{raio} $4$ cm;
  \item o cruzamento deles é o terceiro \omterm{def:g4:geometry:polygon}{vértice} $C$; ligue.
\end{enumerate}
O compasso garante os comprimentos: todo ponto do primeiro arco está
a exatamente $5$ cm de $A$.
\end{method}

\begin{omfigure}
\begin{tikzpicture}[scale=0.95]
  \coordinate (A) at (0,0);
  \coordinate (B) at (6*0.72,0);
  \coordinate (C) at (2.7,2.376);
  \draw[thick] (A) -- (B);
  \draw[thick, omDef] (A) -- (C);
  \draw[thick, omThm] (B) -- (C);
  \draw[omDef, thin] (C) arc[start angle=41.3, delta angle=20,
    radius=3.6];
  \draw[omDef, thin] (C) arc[start angle=41.3, delta angle=-20,
    radius=3.6];
  \draw[omThm, thin] (C) arc[start angle=124.4, delta angle=20,
    radius=2.88];
  \draw[omThm, thin] (C) arc[start angle=124.4, delta angle=-20,
    radius=2.88];
  \fill (A) circle (1.5pt) node[below left, font=\small] {$A$};
  \fill (B) circle (1.5pt) node[below right, font=\small] {$B$};
  \fill (C) circle (1.5pt) node[above, font=\small] {$C$};
  \node[below, font=\small] at (2.16,0) {$6$ cm};
  \node[omDef, font=\small] at (0.85,1.45) {$5$ cm};
  \node[omThm, font=\small] at (3.9,1.35) {$4$ cm};
\end{tikzpicture}

{\small Dois arcos, um cruzamento: o triângulo $6$--$5$--$4$
construído com o compasso. (Se os dois lados curtos juntos fossem
menores que a base, os arcos não se encontrariam --- mais sobre isso
no \cref{ch:g7:triangles}.)}
\end{omfigure}

\begin{example}[Isósceles e equilátero de graça]\label{ex:g5:geometry:special}
Mesma abertura de compasso nos dois arcos $\Rightarrow$ triângulo
isósceles. Abertura igual à base $\Rightarrow$ equilátero. O compasso
não mede: ele \emph{copia} comprimentos, o que é ainda melhor.
\end{example}

\section{Os quadriláteros especiais, por construção}

\begin{method}[Construir a família]\label{met:g5:geometry:quadrilaterals}
\begin{enumerate}
  \item \emph{Retângulo}: dois pares de \omterm{def:g4:geometry:def}{paralelas} que se cruzam em
        \omterm{def:g3:shapes:rightangle}{ângulo reto} --- ou a cadeia do \cref{ex:g5:geometry:chain};
  \item \emph{quadrado}: um retângulo cujos quatro lados usam o mesmo
        comprimento --- transporte-o com o compasso;
  \item \emph{losango}: quatro lados iguais --- trace um lado e
        depois dois arcos de compasso a partir das pontas dele, com a
        mesma abertura, para fixar os outros dois \omterm{def:g4:geometry:polygon}{vértices} (um
        losango são dois triângulos isósceles colados por uma
        diagonal).
\end{enumerate}
Depois de construir, \emph{confira} sempre com as ferramentas:
esquadro nos cantos, compasso nos lados.
\end{method}

\begin{example}[Reconhecer a partir de uma descrição]\label{ex:g5:geometry:recognize}
``Um \omterm{def:g4:geometry:polygon}{quadrilátero} com quatro lados de $4$ cm e um \omterm{def:g3:shapes:rightangle}{ângulo reto}'' ---
construa: os quatro lados iguais forçam um formato de losango, e
empurrar um canto para o \omterm{def:g3:shapes:rightangle}{ângulo reto} endireita todos os outros: o
resultado é um quadrado. Algumas descrições descrevem, sem dizer, uma
forma mais especial do que anunciam. (Por que o \omterm{def:g3:shapes:rightangle}{ângulo reto} se
espalha pelos quatro cantos é explicado no \cref{ch:g7:central}.)
\end{example}

\section{Círculos nas construções}

\begin{example}[O círculo como ferramenta]\label{ex:g5:geometry:circle}
``Encontre todos os pontos a $3$ cm de $A$'': o círculo de centro $A$
e \omterm{def:g3:shapes:shapes}{raio} $3$ cm. ``Encontre todos os pontos a $3$ cm de $A$ \emph{e} a
$4$ cm de $B$'': trace os dois círculos; os cruzamentos deles (dois,
um ou nenhum ponto) são as respostas. Mapas de tesouro e construções
de triângulos usam exatamente essa ideia.
\end{example}

\section{Exercícios}

\begin{exercise}[$\star$]\label{exo:g5:geometry:1}
Trace uma reta $d$ e um ponto $A$ a uns $3$ cm dela. Construa a
perpendicular a $d$ por $A$ e depois a paralela a $d$ por $A$. Marque
os \omterm{def:g3:shapes:rightangle}{ângulos retos}.
\end{exercise}

\begin{exercise}[$\star$]\label{exo:g5:geometry:2}
Construa um retângulo $ABCD$ com $AB = 7$ cm e $BC = 4$ cm, seguindo
o \cref{ex:g5:geometry:chain}. Confira o quarto ângulo com o
esquadro: se a construção for exata, ele é reto automaticamente.
\end{exercise}

\begin{exercise}[$\star$]\label{exo:g5:geometry:3}
Construa um triângulo de lados $7$ cm, $6$ cm e $3$ cm
(\cref{met:g5:geometry:triangle}).
\end{exercise}

\begin{exercise}[$\star$]\label{exo:g5:geometry:4}
Construa um triângulo isósceles de base $5$ cm e lados iguais de
$7$ cm; depois um triângulo equilátero de lado $6$ cm.
\end{exercise}

\begin{exercise}[$\star$]\label{exo:g5:geometry:5}
Construa um quadrado de lado $5$ cm. Que ferramentas garantem os
\omterm{def:g3:shapes:rightangle}{ângulos retos}? Quais garantem os lados iguais?
\end{exercise}

\begin{exercise}[$\star$]\label{exo:g5:geometry:6}
Construa um losango de lado $4$ cm que não seja um quadrado (escolha
você mesmo a abertura do primeiro canto). Confira os quatro lados com
o compasso.
\end{exercise}

\begin{exercise}[$\star$]\label{exo:g5:geometry:7}
Marque dois pontos $A$ e $B$ com $AB = 5$ cm. Construa todos os
pontos que estão a $4$ cm de $A$ e a $3$ cm de $B$. Quantos são?
\end{exercise}

\begin{exercise}[$\star$]\label{exo:g5:geometry:8}
Escreva um programa de construção (passos numerados, pontos com nome)
para um retângulo de $6$ cm $\times$ $2$ cm encimado por um triângulo
isósceles de lados iguais de $4$ cm (o formato de um lápis). Peça a
um colega que o execute.
\end{exercise}

\begin{exercise}[$\star\star$]\label{exo:g5:geometry:9}
Tente construir um triângulo de lados $8$ cm, $3$ cm e $4$ cm. O que
acontece com os arcos? Explique em uma frase por que esse triângulo
não pode existir.
\end{exercise}

\begin{exercise}[$\star\star$]\label{exo:g5:geometry:10}
Construa um losango cujas diagonais meçam $6$ cm e $4$ cm, sabendo
que as diagonais de um losango se cruzam nos seus meios formando
\omterm{def:g3:shapes:rightangle}{ângulo reto}: trace \emph{primeiro} as diagonais e depois ligue as
quatro pontas.
\end{exercise}

\begin{exercise}[$\star\star$]\label{exo:g5:geometry:11}
Uma cabra está amarrada a uma estaca $A$ por uma corda de $4$ m, num
campo plano com uma cerca reta passando a $3$ m de $A$. Faça uma
planta (1 cm por metro): que região a cabra consegue alcançar? Onde a
cerca a limita?
\end{exercise}
